% =====================================================================
\chapter{Extremos libres y condicionados}
\label{tema:5}
% =====================================================================

Entrenar un modelo es minimizar una función de muchas variables, y estimar probabilidades por máxima verosimilitud es maximizar una función con la restricción de que las probabilidades sumen uno. Este tema estudia ambos problemas en dimensión dos y tres: los extremos libres, que se localizan donde se anula el gradiente y se clasifican con la matriz hessiana, y los extremos condicionados, que se resuelven con multiplicadores de Lagrange. La presentación sigue las sesiones~28--31 y~39--40 de 18.02SC~\cite{mit1802}, las §4.7--4.8 de OpenStax~\cite{openstax3}, las §9.9--9.12 y~9.14 de Apostol~\cite{apostol2}, el texto de Marsden y Tromba~\cite{marsden} y las §13.6--13.7 de Strang~\cite{strangcalc}.

\section{Extremos locales y puntos críticos}

\begin{definicion}[Extremos]\label{def:extremos}
Sea $f:D\subseteq\R^n\to\R$ y $\va\in D$. Se dice que $f$ tiene en $\va$ un \emph{máximo local} si existe $\delta>0$ tal que $f(\vx)\le f(\va)$ para todo $\vx\in D$ con $\norm{\vx-\va}<\delta$; un \emph{mínimo local}, con $\ge$. Si la desigualdad vale para todo $\vx\in D$, el extremo es \emph{absoluto} (o global). Un punto $\va$ interior al dominio es \emph{crítico} si $\nabla f(\va)=\vcero$ o si alguna derivada parcial no existe en $\va$.
\end{definicion}

\begin{teorema}[Condición necesaria de Fermat]\label{teo:fermat}
Si $f$ tiene un extremo local en un punto interior $\va$ y sus derivadas parciales existen en $\va$, entonces $\nabla f(\va)=\vcero$.
\end{teorema}

El recíproco es falso: $f(x,y)=x^2-y^2$ tiene $\nabla f(0,0)=\vcero$ y no tiene extremo en el origen, pues crece a lo largo del eje $x$ y decrece a lo largo del eje $y$. Un punto crítico que no es extremo local se llama \emph{punto de silla}.

\section{Clasificación: el criterio del hessiano}

\begin{definicion}[Matriz hessiana]\label{def:hessiana}
Si $f(x,y)$ tiene derivadas segundas, su \emph{matriz hessiana} y su \emph{discriminante} en $\va$ son
\[
H_f(\va)=\begin{pmatrix}f_{xx}&f_{xy}\\f_{yx}&f_{yy}\end{pmatrix}(\va),\qquad
\Delta(\va)=\det H_f(\va)=f_{xx}f_{yy}-f_{xy}^2 ,
\]
donde se ha usado el teorema de Schwarz para escribir $f_{xy}=f_{yx}$.
\end{definicion}

La idea del criterio es la fórmula de Taylor de segundo orden: si $f$ tiene derivadas segundas continuas y $\nabla f(\va)=\vcero$, entonces, con $\mathbf h=(h,k)$,
\begin{equation}\label{eq:t5-taylor}
f(\va+\mathbf h)=f(\va)+\tfrac12\,Q(\mathbf h)+\varepsilon(\mathbf h)\norm{\mathbf h}^2,\qquad Q(h,k)=Ah^2+2Bhk+Ck^2,
\end{equation}
con $A=f_{xx}(\va)$, $B=f_{xy}(\va)$, $C=f_{yy}(\va)$ y $\varepsilon(\mathbf h)\to0$. La demostración de~\eqref{eq:t5-taylor} aplica la fórmula de Taylor de una variable a $t\mapsto f(\va+t\mathbf h)$ (Apostol~\cite[§9.9--9.12]{apostol2}; módulo~4, tema~4). Cerca de un punto crítico, por tanto, el signo de $f(\va+\mathbf h)-f(\va)$ lo decide la \emph{forma cuadrática} $Q$.

\begin{lema}[Signo de una forma cuadrática]\label{lem:forma}
Sea $Q(h,k)=Ah^2+2Bhk+Ck^2$ y $\Delta=AC-B^2$.
\begin{enumerate}
  \item Si $\Delta>0$ y $A>0$, entonces $Q(h,k)>0$ para todo $(h,k)\neq(0,0)$.
  \item Si $\Delta>0$ y $A<0$, entonces $Q(h,k)<0$ para todo $(h,k)\neq(0,0)$.
  \item Si $\Delta<0$, $Q$ toma valores positivos y negativos.
\end{enumerate}
\end{lema}

\begin{teorema}[Criterio del hessiano]\label{teo:hessiano}
Sea $f$ con derivadas segundas continuas cerca de un punto crítico $\va$ con $\nabla f(\va)=\vcero$, y sea $\Delta=\Delta(\va)$.
\begin{enumerate}
  \item Si $\Delta>0$ y $f_{xx}(\va)>0$, $f$ tiene un mínimo local en $\va$.
  \item Si $\Delta>0$ y $f_{xx}(\va)<0$, $f$ tiene un máximo local en $\va$.
  \item Si $\Delta<0$, $\va$ es un punto de silla.
  \item Si $\Delta=0$, el criterio no decide.
\end{enumerate}
\end{teorema}

\addtocounter{definicion}{1}%

\section{Extremos absolutos en regiones cerradas y acotadas}

\begin{teorema}[Weierstrass]\label{teo:weierstrass}
Una función continua sobre un conjunto cerrado y acotado de $\R^n$ alcanza en él un máximo y un mínimo absolutos.
\end{teorema}

La demostración usa la completitud de los números reales y se omite; para una variable es el teorema de los valores extremos del cálculo elemental. En la práctica, para hallar los extremos absolutos de $f$ en una región $R$ cerrada y acotada: (1) se hallan los puntos críticos interiores; (2) se estudia $f$ sobre la frontera, que suele reducirse a funciones de una variable o a un problema con multiplicadores de Lagrange; (3) se comparan todos los valores obtenidos.

\addtocounter{definicion}{1}%

\section{Multiplicadores de Lagrange}

Se busca el máximo o el mínimo de $f(x,y)$ cuando $(x,y)$ está obligado a moverse sobre la curva $g(x,y)=c$. En un extremo, la curva de nivel de $f$ que pasa por el punto no puede cruzar la restricción, porque entonces se podría seguir avanzando sobre $g=c$ hacia valores mayores y menores de $f$: ambas curvas deben ser \emph{tangentes}, y sus gradientes, paralelos (\cref{fig:t5-lagrange}).

\addtocounter{figure}{1}%

\begin{teorema}[Multiplicadores de Lagrange]\label{teo:lagrange}
Sean $f$ y $g$ funciones de $(x,y)$ con derivadas parciales continuas, y sea $\va$ un punto de la curva $g=c$ con $\nabla g(\va)\neq\vcero$. Si $f$, restringida a la curva, tiene un extremo local en $\va$, existe $\lambda\in\R$ (el \emph{multiplicador}) tal que
\[
\nabla f(\va)=\lambda\nabla g(\va).
\]
El mismo enunciado vale para $f$ y $g$ de tres variables y la superficie $g=c$.
\end{teorema}

El método consiste en resolver el sistema formado por $\nabla f=\lambda\nabla g$ y $g=c$ (tres ecuaciones con incógnitas $x,y,\lambda$ en el plano; cuatro en el espacio). Sus soluciones son los \emph{candidatos}; si la restricción es cerrada y acotada, el teorema de Weierstrass garantiza que el mayor valor de $f$ entre ellos es el máximo y el menor, el mínimo.

\addtocounter{definicion}{1}%

\addtocounter{definicion}{1}%

\addtocounter{definicion}{1}%

\addtocounter{definicion}{1}%

